\newpage
% \color{white}
% \pagecolor{black}

\ifdefined\useChapters
\chapter{Desdobrado}
\else
\chapter{}
\fi

Continuo correndo no meio da mata, deixando o vilarejo para trás. O ar entra rasgando a garganta, e cada respiração me lembra o som dos gritos que ficaram ecoando lá atrás. Que massacre. O cheiro de fumaça parece ter grudado em mim. Tenho que correr o máximo que puder, antes que eles me encontrem. Sozinho, não tenho a menor chance contra os quatro.

O chão úmido escorrega sob meus pés, e as folhas secas estalam como se denunciassem meu caminho. A mata fecha em volta, as árvores parecem se inclinar, assistindo à minha fuga. Depois de tanto tempo correndo — nem sei quanto —, consigo despistá-los. O dia já está quase morrendo; o céu, manchado de um vermelho cansado. Estou exausto, com as pernas tremendo. Preciso parar.  
Encosto numa árvore grossa, o tronco frio me apoiando como se também estivesse cansado. “Só um minuto”, penso. Mas o minuto vira uma eternidade.  

É impressionante: uma floresta dessas, e eu nunca havia tentado entrar. O lugar é bonito de um jeito primitivo — uma catedral de verdes e sombras. Apesar de tanta vida, o silêncio é absoluto. Nenhum pássaro, nenhum inseto. Só meu coração batendo alto demais. Talvez seja seguro me encostar aqui, mas, para garantir, rastejo até um tronco caído e me escondo sob ele.  
Deito e tento desacelerar. Minha cabeça é um turbilhão — imagens, vozes, lembranças misturadas. Tento me desligar deles, descansar o corpo, mas algo dentro de mim se recusa a dormir.

Um arrepio percorre meu corpo.  

— Aaaaaah! O que é isso?!

Abro os olhos — e vejo meu corpo ali, imóvel. E eu… aqui, de pé, olhando para ele. De novo.  
Essa sensação já não me assusta tanto quanto antes, mas ainda é estranha: estar e não estar. O corpo parece uma roupa deixada no chão. Bom, já que não tenho a prisão da carne, vou aproveitar. Volitar acima da floresta, ver onde estou.

Subo devagar. É esquisito — não sinto vento, nem peso, nem frio. É como se o ar fosse líquido e eu, uma bolha flutuando dentro dele. Lá em cima, o mundo é silencioso, quase sereno. As copas se juntam como um tapete verde; as sombras formam rios que serpenteiam entre as clareiras.  

Olha só. A cidade lá adiante brilha como se fosse feita de pequenos sóis. As luzes dos carros se cruzam em linhas pulsantes. E, mesmo daqui, consigo ver minha casa. Não fica longe.  

Mas há algo errado.  
O condomínio perto da minha casa… sempre teve quatro prédios, não?  
Agora só vejo três.  

Estranho.  
Um arrepio percorre minha forma sem corpo.

— O que é isso?  
— Como eu vim parar aqui?  
— E esse prédio pegando fogo?

De repente, surge o quarto prédio — o que parecia faltando. Em chamas. Pessoas gritam, acenam das janelas, o brilho das sirenes recorta o céu. O fogo parece vir de um único andar. Não sei o que está acontecendo, mas preciso ver.  
Afinal, estou fora do corpo; não posso me queimar.

Aproximo-me. Quanto mais chego perto, mais o ar vibra, quente, mesmo sem me tocar.  
E então vejo: a menina de cabelo azul. Está ajoelhada, desesperada, labaredas escapando de suas mãos como fitas vivas. Tudo ao redor queima, mas o chão sob ela permanece intacto — um círculo de calma em meio à destruição. O rosto dela está diferente, mais velho, o cabelo mais longo.  

Mas isso é impossível. Há pouco tempo, ela me perseguia na floresta.  
Cabelos não crescem em horas.  
Não… eu não estou apenas mudando de lugar. Estou mudando de tempo.

Uma hipótese absurda começa a se formar. Além de mover-me pelo espaço, talvez eu consiga atravessar o tempo.  
Será parte do desdobramento? Ou essa forma sem corpo obedece a leis que eu nunca imaginei?  
Se for verdade, tempo e espaço não têm mais fronteiras.  

A floresta desaparece, como se alguém apagasse o cenário.  
De repente, estou diante do mesmo prédio, agora intacto. Nenhum fogo. Apenas luzes normais.  
Volito até uma janela acesa e atravesso a parede — sem resistência, como se ela fosse feita de névoa.  
Do outro lado, uma sala simples. E ali, para meu espanto, Crispim está sentado, tomando chá com um rapaz.

Alguém escorrega um bilhete por baixo da porta. O rapaz se move antes mesmo de o papel parar de deslizar, como se soubesse que isso aconteceria. Ele abre a porta, e quem entra é ela — a menina de cabelo azul, viva, desconfiada.  
Eles se encaram com surpresa. É o primeiro encontro.  
Caramba… estou mesmo me movendo no tempo.

A conversa é tensa. Crispim fala com calma, mas o conteúdo é cruel: convence os dois de que é preciso eliminar os que têm habilidades.  
O tom dele não é de fanático; é de alguém que acredita estar sendo racional.  
Pedro e Larissa — os dois parecem perdidos, procurando sentido.  
Eles não têm culpa. São jovens sendo manipulados.  
Meu peito aperta.  
Crispim está por trás de tudo.

Ah! O que é isso agora?  
O chão desaparece, o ar gira.  
Para onde fui? Ou melhor — quando?

Abro os olhos (se é que ainda tenho olhos).  
Estou no mercadinho perto da minha casa.  
O cheiro de farinha e sabão é inconfundível.  

Será que… é o dia em que tudo começou?  
O dia em que o homem me falou para levar o relógio?

Olho ao redor e o vejo — eu mesmo, mais jovem, empilhando as compras no balcão.  
Meu coração dispara.  
Preciso me esconder.  
Pode parecer ridículo, mas depois de ver tantos filmes sobre viagem no tempo, não quero arriscar.  
Fico atrás de uma pilha de caixas, observando meu outro eu.  
Ele pega o último pacote de canela, olha para os lados.  
Mas ninguém aparece.  
Isso é errado.  
Naquele dia, alguém falou comigo.  
Será que mudei alguma coisa?

Se mudei, posso ter quebrado o próprio destino.  
Preciso reparar.

Tomo coragem e crio uma ilusão simples — um disfarce, nada mais.  
Assumo o rosto do ajudante do seu Antônio e me aproximo de mim mesmo.

— Oi, rapaz.  
— Tudo bem com você? Vejo que ainda não encontrou tudo o que procura.

Meu eu do passado sorri, confuso.  
— Já encontrei sim — diz, levantando o pacote de canela.

— Não é disso que estou falando — respondo. — Você sempre foi muito impetuoso e não pretendo te convencer de nada, mas te dou um conselho: hoje à noite, lembre-se de levar seu relógio de bolso.

O eco das palavras vibra no ar.  
A sensação de déjà-vu é quase física — uma vertigem leve, como se o tempo se dobrasse sobre si mesmo.  
Percebo que, na verdade, sempre fui aquele ajudante.  
Tudo estava predestinado.  
Eu não mudei nada; apenas cumpri o que já era inevitável.

E então, sem aviso, o cenário muda novamente.

Estou em outro lugar.  
Reconheço o centro de Arreal — ou o que restou dele.
A prefeitura está destruída, a praça em chamas, carros de ponta-cabeça, postes caídos. O ar é espesso de fumaça e eletricidade.  
Acima de mim, um rapaz flutua, observando o caos como um deus cansado.  
Ao redor, pessoas com poderes se enfrentam, gritam, destroem.  
Algumas lançam fogo; outras erguem carros; outras… nem parecem mais humanas.  
Vejo corpos deformados, olhos brilhando, expressões sem consciência.  
É o futuro.  
E é horrível.

Não posso deixar isso acontecer.  
O poder virou vício.  
As pessoas estão em transe, como marionetes das próprias crenças.  
O mundo inteiro queimando pela vaidade de acreditar demais.

Lembro do discurso de Crispim.  
Ele dizia que queria eliminar alguns para impedir esse futuro.  
Mas isso não é solução.  
É o mesmo erro, só com outra máscara.  
Eu e ele vimos a mesma visão, mas interpretamos de formas opostas.  
Ele quer destruir o perigo.  
Eu quero entender.  
Talvez ainda haja um caminho — uma maneira de impedir o colapso sem repetir o ciclo da destruição.

O vento sopra cinzas no meu rosto invisível.  
Fecho os olhos.  
Preciso agir logo.  
Antes que Crispim, e quem mais ele convencer, transforme o futuro no que acredita ser salvação.
